\section*{Exercises}
\addcontentsline{toc}{section}{Exercises}

The solutions to the exercises that do not ask for code are in
\hyperref[ch:losningar]{appendix~A}. The code is published in the course repository after the
lecture.

\exercise{Find the frame}{Decoding}
\label{ex:2:1}
A node receives the following byte stream:

\begin{console}
32 74 00 FF A5 F7 03 03 25 17 7F 05 10 20 30 02 C2
\end{console}

\begin{enumerate}[a)]
  \item At which index does the frame begin?
  \item Which bytes belong to it, from SOF up to and including CHK?
  \item What is the parser to do with the bytes before it?
  \item What type does the frame have, and which addresses does it travel between?
\end{enumerate}

\exercise{The SOF trap}{Reasoning}
\label{ex:2:2}
The parser is in \code{WaitForSof2} and receives \code{0xA5}.
\begin{enumerate}[a)]
  \item What does the naive implementation do, and what is lost?
  \item What should it do instead?
  \item Construct a byte stream where the difference decides whether a valid frame is found or
    not.
\end{enumerate}

\exercise{A dropped byte}{Decoding}
\label{ex:2:3}
The same frame as in exercise \ref{ex:2:1}, but the byte \code{0x20} is dropped on the way.
\begin{enumerate}[a)]
  \item Which field does the parser read as the checksum?
  \item Will \code{isFrameReady()} become true?
  \item What does \code{extractFrame()} return, and why?
  \item Which state is the parser in afterwards, and what does it take for it to find the next
    frame?
\end{enumerate}

\exercise{The length field}{Reasoning}
\label{ex:2:4}
A parser does not check the length field against the size of its buffer.
\begin{enumerate}[a)]
  \item Construct a byte stream that makes it write past the end of its buffer.
  \item Why is it not enough to check the length in \code{extractFrame()}?
  \item Which general rule follows from that, and where in \kapref{ch:driver} does the same rule
    apply again?
\end{enumerate}
